\section{Testing: From Unit Tests to Chaos}

The project tests at four levels, each catching things the others miss.

\subsection{Level 1: pytest against real services}

The 78 tests in \texttt{tests/} talk to a real Postgres and a real MinIO, not
mocks. Locking behaviour, unique indexes and transactions only exist in the
real database; a mocked database would happily ``pass'' the claim-locking bug.

\paragraph{Isolation.} \texttt{tests/conftest.py} drops and recreates a
\texttt{lab\_test} database and uses a \texttt{lab-test} bucket on every run,
so tests never touch development data:

\begin{lstlisting}[language=Python]
async def _recreate_test_db() -> None:
    conn = await asyncpg.connect(admin_url)
    await conn.execute(f'DROP DATABASE IF EXISTS "{TEST_DB}" WITH (FORCE)')
    await conn.execute(f'CREATE DATABASE "{TEST_DB}"')
    await conn.close()

asyncio.run(_recreate_test_db())
os.environ["DATABASE_URL"] = with_db(BASE_URL, TEST_DB)  # before app import
\end{lstlisting}

The environment variable is set before \texttt{app.config} is imported,
because the settings object reads the environment once, at import time.

\paragraph{Fake agents.} \texttt{harness/client.py} provides a
\texttt{FakeAgent} that speaks the agent protocol through FastAPI's test
client. A test reads like a story:

\begin{lstlisting}[language=Python]
def test_stale_token_rejected(client):
    portal, agent, job_id = isolated(client)
    attempt = agent.claim()
    agent.run_to_complete(attempt)
    bad = agent.complete(attempt, fencing_token=stale_token())
    assert bad["status_code"] == 409
\end{lstlisting}

Each test invents its own GPU class name (\texttt{iso-<random>}), so tests
cannot steal each other's jobs even though they share one database.

\subsection{Level 2: regression tests that fail on the old code}

When a bug is fixed, its test is first run against the \emph{old} code to
prove it catches the bug. In this project that was done for the claim-locking
bug (old code: the second node got nothing), the silent-job bug (old code: hung
past a 15-second limit), the registry-address bug (old code: the job was handed
an unpullable digest) and the Turing change. A test that passes on both old and
new code proves nothing.

\subsection{Level 3: the simulated lab and chaos testing}

Unit tests use fake agents. The simulated lab uses \emph{real} agents in
\emph{real} containers, then breaks things on purpose, an approach called chaos
engineering (Netflix's Chaos Monkey made it famous):

\begin{center}
\begin{tabular}{ll}
\textbf{Scenario} & \textbf{How it breaks things} \\
kill\_agent & \texttt{docker exec pc pkill -f agent.py} mid-job \\
reboot & \texttt{docker restart pc} mid-job \\
network\_drop & \texttt{docker network disconnect} for a full lease \\
cancel\_running & cancel a job while it trains \\
disk\_full & write a tiny number to the fake-disk file \\
desktop\_busy & touch the fake-busy file \\
max\_runtime & a 120 s job with a 5 s limit \\
custom\_env\_via\_registry & build on one PC, pull and run on another \\
self\_update & start a PC on agent 0.0.1 \\
\end{tabular}
\end{center}

Each scenario then asserts the outcome the specification demands, through the
API: the job completed exactly once with one promoted output, it was retried on
another node, the stale attempt was fenced, the cut-off PC rejoined. Short
timers (10 s leases) make the whole exam take about five minutes. It found the
silent-job bug on its first run, after 66 unit tests had passed.

\subsection{Level 4: real hardware}

Simulation cannot show Windows boot behaviour, real GPU drivers or a campus
network. The first real PC found two problems no simulation would: Windows
PowerShell 5.1's stderr handling (see the Windows chapter) and a GPU model missing
from the class list.

Once those were fixed, the first real GPU job ran on that PC: PyTorch inside the
container reported \texttt{cuda\_available True} on an NVIDIA GeForce GTX 1660,
and the output file was promoted. Along the way the agent restarted in the
middle of the job; the server failed that attempt as \texttt{node\_lost} and the
retry on the same PC succeeded. The recovery path the chaos exam rehearses had
just happened for real.

\subsection{Continuous integration}

GitHub Actions runs on every push:

\begin{itemize}
\item \texttt{control}: start Postgres and MinIO with Docker Compose, install
  dependencies, run \texttt{ruff}, run \texttt{alembic check} (are the
  migrations in sync with the models?), run pytest.
\item \texttt{portal}: \texttt{npm ci} and a production \texttt{next build}.
\item \texttt{sim-exam}: nightly and on demand, the full simulated failure exam.
\end{itemize}

CI caught nothing new in this project, which is its job: it makes sure what
passed on a laptop also passes on a clean machine, every time.

\paragraph{The testing pyramid.} Many fast unit tests at the bottom, fewer
slower integration tests in the middle, a few slow end-to-end tests at the top.
This project follows that shape: 78 tests in 16 seconds, 13 chaos scenarios in
5 minutes, one real PC.
